\section{Discussion}

\paragraph{Completion versus validity.}
\ra\ shows that current CLI agents can complete a recognizable research pipeline. They can generate ideas, write code, run experiments, produce figures, draft manuscripts, and review each other. This is a meaningful capability. It is also insufficient. Many failures occur after apparent completion: the experiment is too small, the baseline is weak, the paper reports the wrong setting, or the claimed result cannot be traced to an artifact.

\paragraph{Why final-paper review is insufficient.}
PDF-only review is valuable for measuring apparent paper quality, but it is structurally limited for auto-research. It cannot inspect whether a claimed ablation ran, whether a result file contains the reported number, whether a dataset was synthetic rather than real, or whether a method component exists in code. This explains why \sar\ and \pr\ diverge. The divergence should not be read as a contest between reviewers; it is evidence that paper-only and artifact-grounded protocols evaluate different targets.

\paragraph{Narrative ability exceeds experimental reliability.}
The agents are often better at writing plausible research narratives than at executing reliable scientific studies. Claude Code can synthesize broad, compelling framings, but sometimes overclaims scale or evidence. Codex is more conservative and artifact-faithful, but often produces narrow empirical studies with limited novelty. Kimi Code often proposes formal named systems whose implementation and references do not support the presentation. The common bottleneck is experimental rigor: strong baselines, protocol consistency, statistical care, and faithful result reporting.

\paragraph{Implications for future evaluation.}
Future auto-research evaluations should move beyond final PDFs. At minimum, benchmark releases should preserve code, configs, logs, result artifacts, and review metadata. Review protocols should require reviewers to inspect the relationship between claims and artifacts. Scores should distinguish apparent paper quality from artifact faithfulness and process reliability, because improving one component does not guarantee improvement in the others.
